% Page budget: 1.5--2 pages.
\section{Computational verification}\label{app:numerical-convergence}

\subsection{Exact and floating-point certificates}
\label{app:residual-bounds}

The canonical solver compiles rewards and discounted transition rows from the
registered raw embedded laws in
\path{experiments/configs/unified_canonical_benchmark.toml}.  Its exact mode
uses \texttt{Rational\{BigInt\}} throughout.  Exact rational policy iteration
terminates with Bellman residual zero.  The independently compiled
\texttt{Float64} mode starts from \(v_0=0\), forms
\(v_{n+1}=\mathsf T v_n\), and stops at the first \(n\) satisfying
\begin{equation}\label{eq:numerical-stopping-rule}
  \lVert v_n-v_{n-1}\rVert_\infty\le \tau,
  \qquad \tau=10^{-12},
\end{equation}
subject to the registered cap of 10,000 iterations.  Failure to meet the
tolerance before the cap raises an error rather than emitting a converged
row.

At every iteration the solver separately records the Float64 Bellman residual
\begin{equation}\label{eq:numerical-bellman-residual}
  r_n^{64}:=\lVert\mathsf T_{64} v_n-v_n\rVert_\infty,
\end{equation}
the Float64 a-priori contraction bound
\begin{equation}\label{eq:numerical-apriori-bound}
  \lVert v_n-V_\infty\rVert_\infty
  \le \frac{\beta_{64}^n}{1-\beta_{64}}
       \lVert v_1-v_0\rVert_\infty,
\end{equation}
and the Float64 a-posteriori contraction bound
\begin{equation}\label{eq:numerical-posterior-bound}
  \lVert v_n-V_\infty\rVert_\infty
  \le \frac{r_n^{64}}{1-\beta_{64}}.
\end{equation}
Thus the residual is not substituted for the registered stopping rule.

The solver then applies a separate exact check to the final iterate.  It
converts every binary Float64 value to its exact
\texttt{Rational\{BigInt\}} representation \(\widehat v_{42}\), evaluates
\(\widehat r_{42}=\lVert\mathsf T_{\mathbb Q}\widehat v_{42}-
\widehat v_{42}\rVert_\infty\) with the exact rational Bellman operator, and
forms \(\widehat r_{42}/(1-\beta_{\mathbb Q})\).  This exactly re-evaluated
residual-based certificate is distinct from the Float64 a-posteriori
contraction bound above.

\begin{table}[htbp]
\centering
\small
\begin{tabular}{@{}p{0.43\linewidth}p{0.47\linewidth}@{}}
\toprule
Verification item & Recorded result \\
\midrule
Exact stationary solve & \texttt{Rational\{BigInt\}}; Bellman residual \(0\) \\
Float64 stopping point & \(n=42\); final increment \(5.54223\times10^{-13}\) \\
Final Float64 diagnostics & \(r_{42}^{64}=2.77112\times10^{-13}\); a-priori
contraction bound \(2.52953\times10^{-12}\); a-posteriori contraction bound
\(5.54223\times10^{-13}\) \\
Exactly evaluated Float64 versus exact & maximum absolute error after exact
rationalization \(5.5482\times10^{-13}\), below the exactly re-evaluated
residual-based certificate
\(5.5489\times10^{-13}\) \\
Selected actions & agreement on all six compressed state--belief pairs \\
Outer resource optimization & all eight canonical libraries exhaustively
evaluated for every belief and registered weight schedule; all optimizer ties
retained \\
Envelope bookkeeping & every pairwise crossing candidate recorded separately
from the subset of globally active breakpoints \\
\bottomrule
\end{tabular}
\caption{Compact verification record for the canonical dynamic program and
outer resource optimization.}
\label{tab:canonical-verification-record}
\end{table}

\subsection{Scope of the retained checks}

The dynamic implementation checks 32 raw/compressed embedded transition rows,
144 finite-horizon values, 128 finite-horizon actions, 16 stationary values,
and 16 lifted stationary actions.  It independently evaluates the selected
stationary policy over exact rationals and records every full-duration belief
path, admission atom, and operating-reward block.  The resource implementation
recomputes frontier, closure, burden, productive value, and net value exactly
for each enumerated library.  Numerical solver status is never used as an
exact certificate.

The main paper reports the economically relevant stationary solution.  Online
Supplement~S4 contains the full 42-row convergence history and figure,
horizon-by-horizon values and actions, stationary action values, path and
admission atoms, safe-compression correspondences, capacity paths, penalized
intervals, and all active and inactive pairwise crossing candidates.
